\section{Round-seven reconstruction: symplectic frame-reset transversality, measure-valued trace renewal, and the joint LDP}
\label{sec:r7-b2}

This section replaces the round-six controlling module.  The first-surplus
estimate is reorganized around the youngest separating fork and a transported
orthonormal Jacobi frame.  Its loss is independent of the total ancestral
depth.  Boundary states are finite flux measures, so a stopping contact is not
misidentified with an interior $L^1$ density.  Regularization is performed by
a balance-commuting Markov heat flow and therefore creates no arbitrary
mollification defect.

\subsection{The measure-valued incoming trace hierarchy}

For $k\ge1$ let $\Omega_{k,\varepsilon}$ be the hard-sphere configuration
space and $\partial^-\Omega_{k,\varepsilon}$ its incoming binary-collision
boundary with flux measure
\[
 d\lambda^-_{ij}=| (v_i-v_j)\cdot n_{ij}|\,d\sigma_{ij}\,dV.
\]
A level-$k$ state is a pair $(g_k,\gamma_k^-)$ where $g_k$ is a finite signed
interior measure and $\gamma_k^-$ is a finite signed incoming flux measure.
Atoms are allowed.  On the absolutely continuous part we also retain one
weighted tangential bounded-variation derivative.  With Maxwellian weight
$M_\beta(V)=e^{\beta\sum|v_i|^2}$ define
\[
 \|(g,\gamma)\|_{\alpha,\beta}^{\rm mt}
 =\sum_{k\ge0}\frac{\alpha^k}{k!}
 \left(\int M_\beta\,d|g_k|
 +\int M_\beta\,d|\gamma_k^-|
 +\|D_\tan\gamma_{k,\rm ac}^-\|_{M_\beta}\right).
\]
Specular reflection is an isometry between incoming and outgoing flux
measures.  Free flight sends an outgoing boundary measure into an interior
measure for positive time; it need not regularize at time zero.

Let $U_k(t)$ be free transport with reflection and let $B_{k,k+1}$ be the
creation boundary operator.  The integrated boundary-renewal equation is
\[
 G_k(t)=U_k(t)G_k(0)
 +\int_0^tU_k(t-s)B_{k,k+1}G_{k+1}(s)\,ds,
\]
where $U_k$ acts on the interior--boundary measure pair.

\begin{theorem}[Measure-trace renewal semigroup]
\label{thm:r7-b2-trace}
For sufficiently small activity radius $\alpha$ and every fixed horizon $T$,
the renewal system defines a positive weak-star continuous semigroup
$\mathcal U_\varepsilon(t)$ on the measure-trace hierarchy.  Uniformly for
$0\le t\le T$,
\[
 \|\mathcal U_\varepsilon(t)G\|_{\alpha,\beta}^{\rm mt}
 \le e^{C_Tt}\|G\|_{\alpha,\beta}^{\rm mt}.
\]
The same estimate holds after multiplication of every actual incoming contact
by a bounded source.  A state conditioned at one stopping contact is a boundary
Dirac measure and belongs to this space.  Equivalently, all stopping arguments
may be written as integrated boundary renewal without pointwise conditioning.
\end{theorem}

\begin{proof}
Rescale each collision normal by
$y=(|x_i-x_j|-\varepsilon)/\varepsilon$.  The surface factor and the normal
flux factor then cancel the apparent $\varepsilon$-curvature loss in Green's
formula.  Specular reflection preserves $d\lambda^-$ and the Maxwellian
weight.  Multiple-contact corners have codimension at least two and zero flux
measure; grazing pieces are retained as finite measures and later estimated,
not discarded from the domain.

The creation operator satisfies
\[
 \|B_{k,k+1}G_{k+1}\|_{k}
 \le C(k+1)\|G_{k+1}\|_{k+1}.
\]
The factorial label weight converts this into $C\alpha$ on the full norm.
Picard iteration of the renewal equation therefore converges normally on
$[0,T]$.  Tangential BV is propagated by the reflection derivative on the
nongrazing part; its grazing loss is integrable in flux measure and is absorbed
by the total-variation coordinate.  Bounded contact sources only multiply the
renewal kernel by a fixed factor.  Boundary atoms are admissible from the
start, proving the stopping assertion.
\end{proof}

\subsection{Chronological cells and frame-reset controls}

Fix a collision tree and a chronological cell on which all event orders are
constant.  A first surplus contact joins labels $a,b$.  Let $w$ be their
youngest separating event: either two independent roots, a binary creation
fork, or the last creation event when one label is an ancestor of the other.
Along each exclusive path from $w$ to the surplus contact transport a kinetic
orthonormal transverse frame by the exact linearized free-flight and elastic
scattering cocycle.  At every event perform a QR reset and use the reset
coordinate, rather than the original impact coordinate, as the integration
variable.  Elastic scattering is symplectic in the kinetic metric; the
triangular reset has determinant one when the complementary momentum variable
is transformed simultaneously.

The control block $\Xi_{ab}$ contains at most the following variables per
transverse direction:
\begin{enumerate}[label=(\roman*)]
\item relative root translation for separate-root lineages;
\item the two outgoing impact coordinates and the relative creation time at a
true fork;
\item the last creation impact and a root-velocity shear for an
ancestor--descendant pair.
\end{enumerate}
Common-root translation is not used to control a relative position.

Let $\mathcal R_{ab}$ be the transverse relative-position map at the proposed
surplus contact and
\[
 \mathcal G_{ab}=D_{\Xi_{ab}}\mathcal R_{ab}
                 D_{\Xi_{ab}}\mathcal R_{ab}^{*}.
\]

\begin{theorem}[Depth-independent fork observability]
\label{thm:r7-b2-frame}
On the set where every exclusive event has relative speed at least $\delta$,
incidence angle at least $\delta$, and adjacent event gap at least $\delta$,
there are constants $c,C$ independent of ancestral depth $m$ such that
\[
 \lambda_{\min}\mathcal G_{ab}\ge c\delta^C.
\]
The change from original chronological variables to frame-reset variables has
Jacobian one and preserves the tree integration measure.
\end{theorem}

\begin{proof}
Propagate a terminal covector backwards in the reset frame.  A QR reset makes
the scattering cocycle orthogonal at the level of the chosen transverse
coordinate; its condition number is paid locally when the reset is made and
not multiplied along the whole ancestry.  At a separate-root event,
orthogonality to relative root translations kills the covector.  At a true
fork, orthogonality to the two outgoing impact coordinates kills its two child
components, while the relative creation-time derivative removes the flight
direction.  In the ancestor--descendant case the last impact coordinate kills
the transverse component and the root-velocity shear kills the remaining time
component.  These are the three possible youngest-separation geometries.

Only the local conversion between impact normal, collision time, and the reset
frame uses a denominator.  The nongrazing and event-gap assumptions bound that
conversion by $C\delta^{-C}$, with an exponent depending on dimension but not
on the number of earlier events.  Therefore the adjoint observability estimate
is $\|D\mathcal R_{ab}^{*}\eta\|\ge c\delta^C\|\eta\|$.  Squaring gives the
Gramian bound.  Linearized hard-sphere scattering is symplectic; a symplectic
QR factorization has unit determinant after pairing position and momentum
coordinates, proving measure preservation.
\end{proof}

\subsection{The first-surplus tube and the joint depth sum}

The contact condition has two transverse equations.  Coarea with
Theorem~\ref{thm:r7-b2-frame} gives, on one chronological cell,
\[
 C^m\varepsilon^2\delta^{-C}.
\]
The bad part, where one local speed, angle, or event gap is below $\delta$, has
measure at most $C^m\delta^\beta$ by one-event flux integration.  The exponent
of $\delta$ is independent of $m$ because no composed analytic minor is used.

Choose
\[
 \delta=\varepsilon^\kappa,
 \qquad
 0<\kappa<\frac{2}{C+2C/\beta},
\]
and put
\[
 \alpha_0=\min\{2-\kappa C,\kappa\beta\}>0.
\]

\begin{lemma}[Uniform first-surplus gain]
\label{lem:r7-b2-tube}
For every genealogy of depth $m$, the first-surplus sector is bounded by
\[
 C^m\varepsilon^{\alpha_0}
\]
times the corresponding tree norm.  The exponent $\alpha_0$ is independent of
$m$.
\end{lemma}

\begin{proof}
Add the good coarea bound and the bad one-event flux bound and substitute the
chosen $\delta$.  Theorem~\ref{thm:r7-b2-frame} contains no
$\delta^{-Cm}$ factor, so the stated exponent is uniform.
\end{proof}

Chronological tree expansion contributes $T^m/m!$ and the hierarchy norm
contributes one activity factor per creation.

\begin{theorem}[All-depth cyclic-sector estimate]
\label{thm:r7-b2-cyclic}
On every fixed horizon and bounded source ball,
\[
 \sum_{m\ge0}\frac{(CT\alpha)^m}{m!}
  \,C^m\varepsilon^{\alpha_0}
 \le C_T\varepsilon^{\alpha_0}.
\]
After the first surplus contact the exact reflected microstate is retained and
all later contacts are generated by the measure-trace renewal semigroup.
Therefore the complete cyclic sector, not a collision-deleted surrogate,
vanishes with the displayed power.
\end{theorem}

\begin{proof}
Apply Lemma~\ref{lem:r7-b2-tube} at the first surplus contact and integrate the
future against the outgoing boundary flux measure.  Theorem~\ref{thm:r7-b2-trace}
propagates that measure with norm $e^{C_T(T-t)}$ without point conditioning.
The factorial chronological sum is exponential and independent of
$\varepsilon$, giving the estimate.
\end{proof}

\subsection{Local consistency and fixed-horizon sewing}

Let $\mathcal S_{h,\varepsilon}^H$ be the exact Feynman--Kac block operator for
a source $H$ supported in a time block of length $h$, and let
$\mathcal S_h^H$ be the limiting Boltzmann hierarchy block.  The tree expansion,
Theorem~\ref{thm:r7-b2-cyclic}, and the usual noncyclic Boltzmann--Grad change
of variables give
\[
 \|\mathcal S_{h,\varepsilon}^H-\mathcal S_h^H\|_{R\to R'}
 \le C_R(h\varepsilon^\eta+\varepsilon^{\alpha_0})
\]
on nested measure-trace balls.  The radius grows by at most $e^{Ch}$ per
block.

\begin{theorem}[Stable source sewing]
\label{thm:r7-b2-pressure}
For every fixed $T$ and every bounded particle/contact source in a compact
measure-trace chart, the exact normalized pressure converges locally uniformly,
together with all finite source derivatives, to the sewn Boltzmann pressure.
The total consistency error is
\[
 C_T(\varepsilon^\eta+\varepsilon^{\alpha_0}),
\]
and is uniform in the number of blocks after the mesh is fixed and then sent
to zero.
\end{theorem}

\begin{proof}
Compose the exact and limiting block operators and telescope.  Each preceding
factor is bounded by Theorem~\ref{thm:r7-b2-trace}; each one-block difference
has the displayed consistency bound.  The sum of $T/h$ terms is
$C_T(\varepsilon^\eta+h^{-1}\varepsilon^{\alpha_0})$.  First fix $h$, pass to
the Boltzmann--Grad limit, and then send $h$ to zero.  Normal convergence on a
complex bounded source ball permits Cauchy differentiation.  No reached radius
is hidden: its deterministic bound is $R e^{CT}$.
\end{proof}

\subsection{Balance-preserving regularization}

Let $\mathsf K$ be the Kac collision Markov generator on each fixed
energy--momentum shell and let $\mathsf H$ be velocity-space Ornstein--Uhlenbeck
heat projected onto the collision invariants.  Their product semigroup
$\mathcal R_h=e^{h(\mathsf K+\mathsf H)}$ commutes with the linear balance
operator.  In space and time use convolution after extending by the exact free
transport flow.  Finally mix with a smooth strictly positive feasible baseline
$(f^\circ,\Gamma^\circ)$ by weight $h$.

\begin{lemma}[Exact feasible regularization]
\label{lem:r7-b2-regularize}
Every finite-action feasible pair $(f,\Gamma)$ admits smooth strictly positive
feasible approximants $(f_h,\Gamma_h)$ such that
\[
 (f_h,\Gamma_h)\to(f,\Gamma),
 \qquad
 I(f_h,\Gamma_h)\to I(f,\Gamma).
\]
The balance defect is exactly zero for every $h$; no assertion of an arbitrary
$O(h^M)$ commutator rate is used.
\end{lemma}

\begin{proof}
The Kac and projected heat semigroups preserve mass, momentum, energy, and the
collision incidence relation, hence commute with balance.  Transport-adapted
space--time convolution also commutes with the transport part.  The convex
mixture with the feasible baseline preserves balance and yields strict
positivity.  Markov entropy contraction gives the upper action bound, while
lower semicontinuity gives the reverse liminf.  Approximation of finite
measures by these conservative semigroups proves convergence.
\end{proof}

\subsection{Compactness and the grand-canonical lower bound}

Velocity compactness follows directly from the conserved Maxwellian moment of
the initial law.  Contact compactness follows from the incoming flux estimate
and the entropy inequality
\[
 \Gamma(A)\le
 \frac{I(f,\Gamma)+A_f(A)(e^s-1)}s.
\]
Only bounded indicator sources are used.

For a finite family of bounded tests, Theorem~\ref{thm:r7-b2-pressure} gives a
finite-dimensional differentiable pressure on its regular phase chart.
Strict convexity at that projection follows from the positive microscopic
variance of its exact normalized tilt.  The tilt has a unique projected mean,
and the pressure gap outside every projected neighborhood gives exponential
concentration.  Projecting over an increasing separating family and using
Lemma~\ref{lem:r7-b2-regularize} gives the full lower bound.

\begin{theorem}[Grand-canonical joint actual-contact LDP]
\label{thm:r7-b2-gc}
The empirical density and the actual reflected contact measure satisfy a good
full LDP on every fixed horizon.  Its rate is
\[
 I_0(f_0)+\int_0^T
 \int\ell\left(\frac{d\Gamma}{dA_f}\right)dA_fdt
\]
on balanced pairs and $+\infty$ otherwise.  The lower bound is obtained from
exact finite-volume normalized tilts, finite-dimensional pressure gaps, and
exact feasible regularization; it does not invoke B3 or B4.
\end{theorem}

\begin{proof}
The pressure convergence and compactness above give the projective upper bound.
At a smooth positive pair, solve the finite projected exposing equations and
use the exact finite-volume likelihood.  The pressure gap gives exponential
concentration under the tilt, and change of measure gives the local lower
bound.  Lemma~\ref{lem:r7-b2-regularize} approximates every finite-action pair
without balance error.  Dawson--Gartner and exponential tightness complete the
proof.  Pointwise conjugacy of $e^z-1$ gives the displayed entropy density.
\end{proof}

\subsection{Microcanonical transfer}

Let the B1 coefficient theorem be applied with the correctly extensive shell
and its explicit exponentially small atomic sectors.  Its source-uniform
subexponential coefficient changes no speed-$\mu_\varepsilon$ exponent.

\begin{theorem}[Microcanonical joint LDP]
\label{thm:r7-b2-mc}
Conditioning on the exact particle number and the admissible shrinking
momentum--energy shell restricts the grand-canonical rate of
Theorem~\ref{thm:r7-b2-gc} to the corresponding constraint surface and
subtracts its constrained minimum.  The result holds for the same joint
density--actual-contact topology.
\end{theorem}

\begin{proof}
Apply the B1 source-dependent coefficient to the numerator and denominator of
every finite projected tilt.  The coefficient is subexponential uniformly on
bounded source sets.  The finite-dimensional conditional pressures therefore
converge to the constrained envelopes.  Exponential tightness is inherited
from the conserved shell, and projective upper and lower bounds give the claim.
\end{proof}

\begin{theorem}[Round-seven B2 closure]
\label{thm:r7-b2-main}
The first-surplus loss is independent of ancestral depth; common-root
lineages are controlled at their youngest separating fork; stopping contacts
are legitimate boundary measures; the trace renewal and source sewing are
operator estimates; regularization preserves balance exactly; and the
fixed-horizon grand-canonical and microcanonical joint LDPs follow in the
noncircular order
\[
 \mathrm{B2\mbox{-}GC}\longrightarrow\mathrm{B1}
 \longrightarrow\mathrm{B2\mbox{-}MC}.
\]
\end{theorem}

\begin{proof}
Combine Theorems~\ref{thm:r7-b2-trace},
\ref{thm:r7-b2-frame},
\ref{thm:r7-b2-cyclic},
\ref{thm:r7-b2-pressure},
\ref{thm:r7-b2-gc}, and
\ref{thm:r7-b2-mc}, together with
Lemmas~\ref{lem:r7-b2-tube} and
\ref{lem:r7-b2-regularize}.
\end{proof}
